\documentclass[12pt]{article}
\usepackage{amsmath}
\usepackage{amssymb}
\usepackage{geometry}
\usepackage{listings}
\usepackage{xcolor}
\usepackage{hyperref}
\usepackage{enumitem}
\usepackage{booktabs}
\usepackage{tikz}

\geometry{margin=1in}

\title{MLOps Implementation Plan for Federated Learning DDoS Detection}
\subtitle{Step-by-Step Guide to Operationalize Your FL System}
\author{}
\date{}

\lstset{
    basicstyle=\ttfamily\small,
    keywordstyle=\color{blue},
    stringstyle=\color{red},
    commentstyle=\color{green},
    breaklines=true,
    frame=single
}

\begin{document}

\maketitle

\section{Overview}

This document provides a comprehensive, step-by-step plan to implement MLOps for your Federated Learning DDoS detection system. The plan is divided into phases with specific tasks, code examples, and success criteria.

\section{Current State Assessment}

\subsection{What You Have}
\begin{itemize}
    \item ✅ Federated Learning infrastructure (Flower)
    \item ✅ Basic monitoring dashboard
    \item ✅ Model training and evaluation scripts
    \item ✅ Docker containerization
    \item ✅ Multiple model architectures (MLP, LSTM, CNN1D, CNN\_LSTM)
\end{itemize}

\subsection{What's Missing}
\begin{itemize}
    \item ❌ Experiment tracking system
    \item ❌ Model versioning and registry
    \item ❌ Automated CI/CD pipelines
    \item ❌ Production deployment
    \item ❌ Automated retraining
    \item ❌ Data versioning
    \item ❌ Model monitoring in production
\end{itemize}

\section{Implementation Phases}

\subsection{Phase 1: Foundation - Experiment Tracking \& Versioning (Week 1-2)}

\subsubsection{Step 1.1: Install MLOps Tools}

\textbf{Tasks:}
\begin{enumerate}
    \item Install MLflow for experiment tracking
    \item Install DVC for data versioning (optional)
    \item Set up MLflow tracking server
\end{enumerate}

\textbf{Commands:}
\begin{lstlisting}[language=bash]
# Install MLflow
pip install mlflow

# Install DVC (optional, for data versioning)
pip install dvc dvc-s3  # or dvc-gdrive

# Start MLflow tracking server
mlflow server --backend-store-uri sqlite:///mlflow.db \
              --default-artifact-root ./mlruns \
              --host 0.0.0.0 --port 5000
\end{lstlisting}

\textbf{Update requirements.txt:}
\begin{lstlisting}[language=text]
mlflow>=2.8.0
dvc>=3.0.0  # Optional
\end{lstlisting}

\subsubsection{Step 1.2: Integrate MLflow into FL Server}

\textbf{File:} `flower_server_metrics.py`

\textbf{Changes:}
\begin{lstlisting}[language=Python]
import mlflow
import mlflow.keras

class MetricsFedAvg(FedAvg):
    def __init__(self, ...):
        super().__init__(...)
        # Initialize MLflow
        mlflow.set_tracking_uri("http://localhost:5000")
        mlflow.set_experiment("DDoS_Detection_FL")
        self.mlflow_run = None
    
    def configure_fit(self, server_round, parameters, client_manager):
        # Start MLflow run for this FL training session
        if server_round == 1:
            self.mlflow_run = mlflow.start_run(
                run_name=f"{self.model_type}_FL_{datetime.now().strftime('%Y%m%d_%H%M%S')}"
            )
            # Log FL configuration
            mlflow.log_params({
                "model_type": self.model_type,
                "num_rounds": self.num_rounds,
                "min_clients": self.min_fit_clients,
                "fraction_fit": self.fraction_fit,
                "aggregation": "FedAvg"
            })
        
        # Log round start
        mlflow.log_metric("round", server_round, step=server_round)
        
        return super().configure_fit(server_round, parameters, client_manager)
    
    def aggregate_fit(self, rnd, results, failures):
        aggregated_weights, aggregated_metrics = super().aggregate_fit(rnd, results, failures)
        
        # Log round metrics
        accuracy = aggregated_metrics.get("accuracy", 0.0)
        loss = aggregated_metrics.get("loss", 0.0)
        
        mlflow.log_metrics({
            "round_accuracy": float(accuracy),
            "round_loss": float(loss),
            "bytes_sent": self.round_metrics["bytes_sent"],
            "bytes_received": self.round_metrics["bytes_received"],
            "round_time": round_time
        }, step=rnd)
        
        # Save model checkpoint for this round
        if self.mlflow_run:
            model = self._create_model_architecture()
            # Set aggregated weights
            # ... (weight assignment code) ...
            mlflow.keras.log_model(
                model, 
                f"model_round_{rnd}",
                registered_model_name=f"{self.model_type}_FL"
            )
        
        # End run after final round
        if rnd >= self.num_rounds and self.mlflow_run:
            mlflow.end_run()
        
        return aggregated_weights, aggregated_metrics
\end{lstlisting}

\textbf{Success Criteria:}
\begin{itemize}
    \item MLflow server running and accessible
    \item Each FL training session creates MLflow run
    \item Round metrics logged to MLflow
    \item Models saved as MLflow artifacts
\end{itemize}

\subsubsection{Step 1.3: Implement Model Versioning}

\textbf{File:} Create `mlops/model_registry.py`

\textbf{Code:}
\begin{lstlisting}[language=Python]
import os
import mlflow
from datetime import datetime
from tensorflow.keras.models import load_model

class FLModelRegistry:
    def __init__(self, tracking_uri="http://localhost:5000"):
        mlflow.set_tracking_uri(tracking_uri)
        self.client = mlflow.tracking.MlflowClient()
    
    def register_model_version(self, model_path, model_type, round_num, 
                               accuracy, loss, metadata=None):
        """Register a new model version after FL round"""
        timestamp = datetime.now().strftime("%Y%m%d_%H%M%S")
        version_name = f"{model_type}_FL_Round{round_num}_Acc{accuracy:.3f}_{timestamp}"
        
        # Register in MLflow
        model_name = f"{model_type}_FL"
        
        # Log model
        run = mlflow.start_run(run_name=version_name)
        mlflow.log_params({
            "model_type": model_type,
            "round": round_num,
            "version": version_name
        })
        mlflow.log_metrics({
            "accuracy": accuracy,
            "loss": loss
        })
        mlflow.keras.log_model(
            load_model(model_path),
            "model",
            registered_model_name=model_name
        )
        mlflow.end_run()
        
        return version_name
    
    def get_latest_model(self, model_type):
        """Get latest registered model version"""
        model_name = f"{model_type}_FL"
        latest_version = self.client.get_latest_versions(
            model_name, stages=["None"]
        )
        if latest_version:
            return latest_version[0]
        return None
    
    def promote_to_production(self, model_type, version):
        """Promote model version to production"""
        model_name = f"{model_type}_FL"
        self.client.transition_model_version_stage(
            name=model_name,
            version=version,
            stage="Production"
        )
\end{lstlisting}

\textbf{Integration:}
\begin{lstlisting}[language=Python]
# In flower_server_metrics.py, after aggregation:
from mlops.model_registry import FLModelRegistry

registry = FLModelRegistry()
version_name = registry.register_model_version(
    model_path=model_path,
    model_type=self.model_type,
    round_num=rnd,
    accuracy=accuracy,
    loss=loss
)
\end{lstlisting}

\subsubsection{Step 1.4: Enhanced Logging}

\textbf{File:} Create `mlops/fl_logger.py`

\textbf{Code:}
\begin{lstlisting}[language=Python]
import logging
import json
from datetime import datetime
from pathlib import Path

class FLLogger:
    def __init__(self, log_dir="logs/fl_mlops"):
        Path(log_dir).mkdir(parents=True, exist_ok=True)
        self.log_dir = log_dir
        
        # Structured logging
        self.logger = logging.getLogger("FL_MLOps")
        self.logger.setLevel(logging.INFO)
        
        # File handler
        log_file = Path(log_dir) / f"fl_{datetime.now().strftime('%Y%m%d')}.log"
        file_handler = logging.FileHandler(log_file)
        file_handler.setFormatter(
            logging.Formatter('%(asctime)s - %(name)s - %(levelname)s - %(message)s')
        )
        self.logger.addHandler(file_handler)
    
    def log_round(self, round_num, metrics):
        """Log structured round metrics"""
        log_entry = {
            "timestamp": datetime.now().isoformat(),
            "round": round_num,
            "metrics": metrics
        }
        
        # Write to JSON log
        json_log = Path(self.log_dir) / f"round_{round_num}.json"
        with open(json_log, 'w') as f:
            json.dump(log_entry, f, indent=2)
        
        # Also log to standard logger
        self.logger.info(f"Round {round_num}: {json.dumps(metrics)}")
    
    def log_model_version(self, model_type, version, metrics):
        """Log model version information"""
        log_entry = {
            "timestamp": datetime.now().isoformat(),
            "model_type": model_type,
            "version": version,
            "metrics": metrics
        }
        
        json_log = Path(self.log_dir) / f"model_{model_type}_{version}.json"
        with open(json_log, 'w') as f:
            json.dump(log_entry, f, indent=2)
\end{lstlisting}

\subsection{Phase 2: Data Management \& Validation (Week 2-3)}

\subsubsection{Step 2.1: Data Versioning with DVC}

\textbf{Setup:}
\begin{lstlisting}[language=bash]
# Initialize DVC
dvc init

# Add dataset to DVC
dvc add dataset_sdn.csv

# Commit to git
git add dataset_sdn.csv.dvc .gitignore
git commit -m "Add dataset to DVC"
\end{lstlisting}

\textbf{Create `.dvc/config`:}
\begin{lstlisting}[language=ini]
['remote "local"']
url = ./dvc_storage
['remote "s3"']
url = s3://your-bucket/datasets
\end{lstlisting}

\subsubsection{Step 2.2: Data Validation Pipeline}

\textbf{File:} Create `mlops/data_validation.py`

\textbf{Code:}
\begin{lstlisting}[language=Python]
import pandas as pd
import numpy as np
from typing import Dict, Tuple
import great_expectations as ge  # pip install great-expectations

class DataValidator:
    def __init__(self):
        self.expectations = []
    
    def validate_dataset(self, df: pd.DataFrame) -> Tuple[bool, Dict]:
        """Validate dataset before FL training"""
        ge_df = ge.from_pandas(df)
        
        # Define expectations
        expectations = {
            "row_count": len(df) > 0,
            "no_null_labels": df['label'].notna().all(),
            "label_binary": df['label'].isin([0, 1]).all(),
            "features_range": self._check_feature_ranges(df),
            "class_balance": self._check_class_balance(df)
        }
        
        results = {}
        all_passed = True
        
        for name, expectation in expectations.items():
            if isinstance(expectation, bool):
                results[name] = expectation
                if not expectation:
                    all_passed = False
            else:
                results[name] = expectation
        
        return all_passed, results
    
    def _check_feature_ranges(self, df):
        """Check if features are within expected ranges"""
        numeric_cols = df.select_dtypes(include=[np.number]).columns
        numeric_cols = [c for c in numeric_cols if c != 'label']
        
        issues = []
        for col in numeric_cols:
            if df[col].min() < 0 and col not in ['port_no']:  # Some cols can be negative
                if 'count' in col.lower() or 'byte' in col.lower():
                    issues.append(f"{col} has negative values")
        
        return len(issues) == 0, issues
    
    def _check_class_balance(self, df):
        """Check class distribution"""
        class_counts = df['label'].value_counts()
        balance_ratio = class_counts.min() / class_counts.max()
        
        # Warn if severely imbalanced (< 0.3)
        is_balanced = balance_ratio >= 0.3
        return is_balanced, {
            "ratio": balance_ratio,
            "class_0": int(class_counts[0]),
            "class_1": int(class_counts[1])
        }
\end{lstlisting}

\textbf{Integration:}
\begin{lstlisting}[language=Python]
# In flower_worker.py, _load_data method:
from mlops.data_validation import DataValidator

def _load_data(self):
    df = pd.read_csv("dataset_sdn.csv")
    # ... preprocessing ...
    
    # Validate data
    validator = DataValidator()
    is_valid, validation_results = validator.validate_dataset(df)
    
    if not is_valid:
        logger.warning(f"Data validation issues: {validation_results}")
        # Could raise exception or log warning
    
    # Continue with data loading...
\end{lstlisting}

\subsubsection{Step 2.3: Data Quality Monitoring}

\textbf{File:} Create `mlops/data_quality.py`

\textbf{Code:}
\begin{lstlisting}[language=Python]
import pandas as pd
import numpy as np
from datetime import datetime

class DataQualityMonitor:
    def __init__(self):
        self.baseline_stats = None
    
    def compute_baseline(self, df: pd.DataFrame):
        """Compute baseline statistics for data quality checks"""
        self.baseline_stats = {
            "mean": df.select_dtypes(include=[np.number]).mean().to_dict(),
            "std": df.select_dtypes(include=[np.number]).std().to_dict(),
            "min": df.select_dtypes(include=[np.number]).min().to_dict(),
            "max": df.select_dtypes(include=[np.number]).max().to_dict(),
            "null_counts": df.isnull().sum().to_dict(),
            "class_distribution": df['label'].value_counts().to_dict()
        }
    
    def detect_drift(self, new_df: pd.DataFrame, threshold=0.1) -> Dict:
        """Detect data drift compared to baseline"""
        if self.baseline_stats is None:
            raise ValueError("Baseline not computed")
        
        drift_detected = {}
        
        # Check feature distributions
        for col in new_df.select_dtypes(include=[np.number]).columns:
            if col in self.baseline_stats['mean']:
                baseline_mean = self.baseline_stats['mean'][col]
                new_mean = new_df[col].mean()
                
                drift_ratio = abs(new_mean - baseline_mean) / (baseline_mean + 1e-6)
                if drift_ratio > threshold:
                    drift_detected[col] = {
                        "baseline_mean": baseline_mean,
                        "new_mean": new_mean,
                        "drift_ratio": drift_ratio
                    }
        
        # Check class distribution
        new_class_dist = new_df['label'].value_counts().to_dict()
        baseline_dist = self.baseline_stats['class_distribution']
        
        for label in [0, 1]:
            baseline_ratio = baseline_dist.get(label, 0) / sum(baseline_dist.values())
            new_ratio = new_class_dist.get(label, 0) / sum(new_class_dist.values())
            
            if abs(new_ratio - baseline_ratio) > threshold:
                drift_detected[f"class_{label}"] = {
                    "baseline_ratio": baseline_ratio,
                    "new_ratio": new_ratio
                }
        
        return drift_detected
\end{lstlisting}

\subsection{Phase 3: CI/CD Pipeline (Week 3-4)}

\subsubsection{Step 3.1: GitHub Actions Workflow}

\textbf{File:} Create `.github/workflows/fl_training.yml`

\textbf{Code:}
\begin{lstlisting}[language=YAML]
name: Federated Learning Training Pipeline

on:
  schedule:
    # Weekly retraining every Sunday at 2 AM
    - cron: '0 2 * * 0'
  workflow_dispatch:
    inputs:
      model_type:
        description: 'Model type to train'
        required: true
        default: 'MLPv2'
        type: choice
        options:
          - MLPv2
          - LSTM
          - CNN1D
          - CNN_LSTM
      num_rounds:
        description: 'Number of FL rounds'
        required: false
        default: 5
        type: number

jobs:
  validate-data:
    runs-on: ubuntu-latest
    steps:
      - uses: actions/checkout@v3
      
      - name: Set up Python
        uses: actions/setup-python@v4
        with:
          python-version: '3.9'
      
      - name: Install dependencies
        run: |
          pip install -r requirements.txt
          pip install great-expectations
      
      - name: Validate dataset
        run: |
          python mlops/validate_data.py
      
      - name: Check data quality
        run: |
          python mlops/check_data_quality.py

  fl-training:
    needs: validate-data
    runs-on: ubuntu-latest
    strategy:
      matrix:
        model_type: [MLPv2, LSTM, CNN1D, CNN_LSTM]
    
    steps:
      - uses: actions/checkout@v3
      
      - name: Set up Python
        uses: actions/setup-python@v4
        with:
          python-version: '3.9'
      
      - name: Install dependencies
        run: |
          pip install -r requirements.txt
          pip install mlflow
      
      - name: Start MLflow server
        run: |
          mlflow server --backend-store-uri sqlite:///mlflow.db \
                        --default-artifact-root ./mlruns \
                        --host 0.0.0.0 --port 5000 &
        continue-on-error: true
      
      - name: Start FL Server
        run: |
          docker-compose up -d flower-server-${{ matrix.model_type.lower() }}
        env:
          FL_MODEL_TYPE: ${{ matrix.model_type }}
      
      - name: Start FL Workers
        run: |
          docker-compose up -d flower-worker-1 flower-worker-2
      
      - name: Wait for FL Training
        run: |
          python mlops/wait_for_fl_completion.py \
            --model-type ${{ matrix.model_type }} \
            --max-wait 3600
      
      - name: Evaluate Model
        run: |
          python evaluate.py --model-name ${{ matrix.model_type }}_FL.h5
      
      - name: Register Model in MLflow
        run: |
          python mlops/register_model.py \
            --model-path models/${{ matrix.model_type }}_FL.h5 \
            --model-type ${{ matrix.model_type }}
      
      - name: Compare with Previous Model
        run: |
          python mlops/compare_models.py \
            --new-model ${{ matrix.model_type }}_FL.h5 \
            --model-type ${{ matrix.model_type }}
      
      - name: Deploy if Improved
        if: success()
        run: |
          python mlops/deploy_if_improved.py \
            --model-type ${{ matrix.model_type }}

  notify:
    needs: [fl-training]
    runs-on: ubuntu-latest
    if: always()
    steps:
      - name: Notify Results
        run: |
          echo "FL Training completed for all models"
          # Add Slack/Email notification here
\end{lstlisting}

\subsubsection{Step 3.2: Pre-commit Hooks}

\textbf{File:} Create `.pre-commit-config.yaml`

\textbf{Code:}
\begin{lstlisting}[language=YAML]
repos:
  - repo: https://github.com/pre-commit/pre-commit-hooks
    rev: v4.4.0
    hooks:
      - id: trailing-whitespace
      - id: end-of-file-fixer
      - id: check-yaml
      - id: check-added-large-files
      - id: check-json
  
  - repo: https://github.com/psf/black
    rev: 23.3.0
    hooks:
      - id: black
        language_version: python3.9
  
  - repo: https://github.com/pycqa/flake8
    rev: 6.0.0
    hooks:
      - id: flake8
        args: [--max-line-length=100]
\end{lstlisting}

\textbf{Setup:}
\begin{lstlisting}[language=bash]
pip install pre-commit
pre-commit install
\end{lstlisting}

\subsubsection{Step 3.3: Automated Testing}

\textbf{File:} Create `tests/test_fl_pipeline.py`

\textbf{Code:}
\begin{lstlisting}[language=Python]
import pytest
import numpy as np
from flower_worker import FlowerWorker
from mlops.data_validation import DataValidator

class TestFLPipeline:
    def test_data_validation(self):
        """Test data validation before FL training"""
        validator = DataValidator()
        # Create test dataframe
        df = pd.DataFrame({
            'feature1': [1, 2, 3],
            'feature2': [4, 5, 6],
            'label': [0, 1, 0]
        })
        
        is_valid, results = validator.validate_dataset(df)
        assert is_valid == True
    
    def test_model_creation(self):
        """Test model creation for each type"""
        model_types = ['MLPv2', 'CNN1D', 'LSTM', 'CNN_LSTM']
        
        for model_type in model_types:
            worker = FlowerWorker(
                worker_id="test_worker",
                model_type=model_type,
                data_partition=0.1
            )
            model = worker._create_model()
            assert model is not None
            assert model.count_params() > 0
    
    def test_weight_aggregation(self):
        """Test FedAvg weight aggregation"""
        # Simulate weights from 2 clients
        weights_client1 = [np.array([1.0, 2.0]), np.array([3.0])]
        weights_client2 = [np.array([2.0, 3.0]), np.array([4.0])]
        
        # Simulate FedAvg: average weights
        aggregated = [
            (w1 + w2) / 2 
            for w1, w2 in zip(weights_client1, weights_client2)
        ]
        
        assert len(aggregated) == 2
        assert np.allclose(aggregated[0], [1.5, 2.5])
\end{lstlisting}

\textbf{Run tests:}
\begin{lstlisting}[language=bash]
pip install pytest pytest-cov
pytest tests/ -v --cov=.
\end{lstlisting}

\subsection{Phase 4: Model Deployment (Week 4-5)}

\subsubsection{Step 4.1: Create Model Serving API}

\textbf{File:} Create `mlops/model_serving.py`

\textbf{Code:}
\begin{lstlisting}[language=Python]
from fastapi import FastAPI, HTTPException
from pydantic import BaseModel
import numpy as np
import pandas as pd
from tensorflow.keras.models import load_model
from mlops.model_registry import FLModelRegistry
import mlflow

app = FastAPI(title="DDoS Detection API")

# Load model registry
registry = FLModelRegistry()

# Cache for loaded models
model_cache = {}

class PredictionRequest(BaseModel):
    features: list
    model_type: str = "MLPv2"

class PredictionResponse(BaseModel):
    prediction: int
    probability: float
    model_version: str

@app.on_event("startup")
async def load_models():
    """Load all models on startup"""
    model_types = ['MLPv2', 'LSTM', 'CNN1D', 'CNN_LSTM']
    
    for model_type in model_types:
        try:
            # Get latest production model
            latest = registry.get_latest_model(model_type)
            if latest:
                model_path = mlflow.artifacts.download_artifacts(
                    run_id=latest.run_id,
                    artifact_path="model"
                )
                model_cache[model_type] = load_model(model_path)
        except Exception as e:
            print(f"Warning: Could not load {model_type}: {e}")

@app.get("/health")
async def health_check():
    """Health check endpoint"""
    return {
        "status": "healthy",
        "models_loaded": list(model_cache.keys())
    }

@app.post("/predict", response_model=PredictionResponse)
async def predict(request: PredictionRequest):
    """Predict DDoS attack"""
    if request.model_type not in model_cache:
        raise HTTPException(
            status_code=404,
            detail=f"Model {request.model_type} not available"
        )
    
    model = model_cache[request.model_type]
    
    # Prepare input
    features = np.array([request.features])
    
    # Reshape for CNN/LSTM models
    if request.model_type in ['CNN1D', 'LSTM', 'CNN_LSTM']:
        features = np.expand_dims(features, axis=2)
    
    # Predict
    prediction_proba = model.predict(features, verbose=0)[0]
    prediction = int(np.argmax(prediction_proba))
    probability = float(prediction_proba[prediction])
    
    # Get model version
    latest = registry.get_latest_model(request.model_type)
    model_version = latest.version if latest else "unknown"
    
    return PredictionResponse(
        prediction=prediction,
        probability=probability,
        model_version=model_version
    )

@app.get("/models")
async def list_models():
    """List available models"""
    return {
        "available_models": list(model_cache.keys()),
        "model_info": {
            model_type: {
                "parameters": model.count_params(),
                "input_shape": str(model.input_shape),
                "output_shape": str(model.output_shape)
            }
            for model_type, model in model_cache.items()
        }
    }

if __name__ == "__main__":
    import uvicorn
    uvicorn.run(app, host="0.0.0.0", port=8000)
\end{lstlisting}

\subsubsection{Step 4.2: Dockerize Model Serving}

\textbf{File:} Create `Dockerfile.serving`

\textbf{Code:}
\begin{lstlisting}[language=Dockerfile]
FROM python:3.9-slim

WORKDIR /app

# Install dependencies
COPY requirements-serving.txt .
RUN pip install --no-cache-dir -r requirements-serving.txt

# Copy application
COPY mlops/model_serving.py .
COPY mlops/model_registry.py ./mlops/

# Expose port
EXPOSE 8000

# Run API server
CMD ["python", "-m", "uvicorn", "mlops.model_serving:app", \
     "--host", "0.0.0.0", "--port", "8000"]
\end{lstlisting}

\textbf{File:} Create `requirements-serving.txt`

\textbf{Code:}
\begin{lstlisting}[language=text]
fastapi>=0.104.0
uvicorn>=0.24.0
tensorflow>=2.13.0
mlflow>=2.8.0
pydantic>=2.0.0
numpy>=1.24.0
pandas>=2.0.0
\end{lstlisting}

\subsubsection{Step 4.3: Add to docker-compose.yml}

\textbf{Add to docker-compose.yml:}
\begin{lstlisting}[language=YAML]
  # Model Serving API
  model-serving:
    build:
      context: .
      dockerfile: Dockerfile.serving
    container_name: model-serving-api
    ports:
      - "8000:8000"
    environment:
      - MLFLOW_TRACKING_URI=http://mlflow-server:5000
    volumes:
      - ./models:/app/models
      - ./mlruns:/app/mlruns
    networks:
      - flower-network
    depends_on:
      - mlflow-server

  # MLflow Tracking Server
  mlflow-server:
    image: python:3.9-slim
    container_name: mlflow-server
    ports:
      - "5000:5000"
    volumes:
      - ./mlruns:/mlruns
      - ./mlflow.db:/mlflow.db
    command: >
      bash -c "pip install mlflow &&
               mlflow server --backend-store-uri sqlite:///mlflow.db
                            --default-artifact-root ./mlruns
                            --host 0.0.0.0 --port 5000"
    networks:
      - flower-network
\end{lstlisting}

\subsection{Phase 5: Monitoring \& Alerting (Week 5-6)}

\subsubsection{Step 5.1: Production Monitoring}

\textbf{File:} Create `mlops/production_monitor.py`

\textbf{Code:}
\begin{lstlisting}[language=Python]
import time
import requests
import numpy as np
from datetime import datetime
from typing import Dict, List
import mlflow

class ProductionMonitor:
    def __init__(self, api_url="http://localhost:8000"):
        self.api_url = api_url
        self.predictions_history = []
        self.performance_metrics = {
            "total_predictions": 0,
            "ddos_detected": 0,
            "avg_latency": 0.0,
            "error_rate": 0.0
        }
    
    def log_prediction(self, features: List, prediction: int, 
                      probability: float, latency: float):
        """Log prediction for monitoring"""
        self.predictions_history.append({
            "timestamp": datetime.now().isoformat(),
            "features": features,
            "prediction": prediction,
            "probability": probability,
            "latency": latency
        })
        
        # Update metrics
        self.performance_metrics["total_predictions"] += 1
        if prediction == 1:
            self.performance_metrics["ddos_detected"] += 1
        
        # Update average latency
        total = self.performance_metrics["total_predictions"]
        current_avg = self.performance_metrics["avg_latency"]
        self.performance_metrics["avg_latency"] = (
            (current_avg * (total - 1) + latency) / total
        )
        
        # Log to MLflow
        mlflow.log_metrics({
            "predictions_per_minute": total / (time.time() - self.start_time) * 60,
            "ddos_detection_rate": self.performance_metrics["ddos_detected"] / total,
            "avg_latency_ms": self.performance_metrics["avg_latency"] * 1000
        })
    
    def check_performance_degradation(self, threshold_accuracy=0.85):
        """Check if model performance degraded"""
        # This would require ground truth labels
        # In production, you'd compare predictions with actual outcomes
        # For now, monitor prediction confidence
        
        if len(self.predictions_history) < 100:
            return False
        
        recent_predictions = self.predictions_history[-100:]
        avg_confidence = np.mean([p["probability"] for p in recent_predictions])
        
        if avg_confidence < threshold_accuracy:
            return True
        return False
    
    def detect_data_drift(self, baseline_features, new_features, threshold=0.1):
        """Detect if input data distribution changed"""
        from mlops.data_quality import DataQualityMonitor
        
        monitor = DataQualityMonitor()
        monitor.compute_baseline(baseline_features)
        drift = monitor.detect_drift(new_features, threshold)
        
        return len(drift) > 0, drift
\end{lstlisting}

\subsubsection{Step 5.2: Alerting System}

\textbf{File:} Create `mlops/alerting.py`

\textbf{Code:}
\begin{lstlisting}[language=Python]
import smtplib
from email.mime.text import MIMEText
from typing import Dict, List

class AlertManager:
    def __init__(self, config: Dict):
        self.config = config
        self.alert_history = []
    
    def send_alert(self, alert_type: str, message: str, severity: str = "warning"):
        """Send alert via configured channels"""
        alert = {
            "timestamp": datetime.now().isoformat(),
            "type": alert_type,
            "message": message,
            "severity": severity
        }
        
        self.alert_history.append(alert)
        
        # Email alert
        if self.config.get("email_enabled", False):
            self._send_email(alert)
        
        # Slack alert
        if self.config.get("slack_enabled", False):
            self._send_slack(alert)
        
        # Log alert
        print(f"[{severity.upper()}] {alert_type}: {message}")
    
    def _send_email(self, alert: Dict):
        """Send email alert"""
        # Implementation for email alerts
        pass
    
    def _send_slack(self, alert: Dict):
        """Send Slack alert"""
        # Implementation for Slack webhook
        pass
    
    def check_and_alert(self, monitor: ProductionMonitor):
        """Check monitoring metrics and send alerts if needed"""
        # Performance degradation
        if monitor.check_performance_degradation():
            self.send_alert(
                "performance_degradation",
                "Model performance below threshold",
                "critical"
            )
        
        # High error rate
        if monitor.performance_metrics["error_rate"] > 0.05:
            self.send_alert(
                "high_error_rate",
                f"Error rate: {monitor.performance_metrics['error_rate']:.2%}",
                "warning"
            )
        
        # High latency
        if monitor.performance_metrics["avg_latency"] > 1.0:  # > 1 second
            self.send_alert(
                "high_latency",
                f"Average latency: {monitor.performance_metrics['avg_latency']:.2f}s",
                "warning"
            )
\end{lstlisting}

\subsubsection{Step 5.3: Prometheus Metrics Export}

\textbf{File:} Create `mlops/metrics_exporter.py`

\textbf{Code:}
\begin{lstlisting}[language=Python]
from prometheus_client import Counter, Histogram, Gauge, start_http_server

# Define metrics
predictions_total = Counter(
    'ddos_predictions_total',
    'Total number of predictions',
    ['model_type', 'prediction']
)

prediction_latency = Histogram(
    'ddos_prediction_latency_seconds',
    'Prediction latency in seconds',
    ['model_type']
)

model_accuracy = Gauge(
    'ddos_model_accuracy',
    'Current model accuracy',
    ['model_type', 'version']
)

def export_metrics(port=9090):
    """Start Prometheus metrics server"""
    start_http_server(port)
    print(f"Prometheus metrics server started on port {port}")

# In model_serving.py, add metrics:
@app.post("/predict")
async def predict(request: PredictionRequest):
    start_time = time.time()
    
    # ... prediction logic ...
    
    # Export metrics
    predictions_total.labels(
        model_type=request.model_type,
        prediction=prediction
    ).inc()
    
    prediction_latency.labels(
        model_type=request.model_type
    ).observe(time.time() - start_time)
\end{lstlisting}

\subsection{Phase 6: Automated Retraining (Week 6-7)}

\subsubsection{Step 6.1: Retraining Trigger Logic}

\textbf{File:} Create `mlops/retraining_trigger.py`

\textbf{Code:}
\begin{lstlisting}[language=Python]
from datetime import datetime, timedelta
from mlops.production_monitor import ProductionMonitor
from mlops.model_registry import FLModelRegistry
import mlflow

class RetrainingTrigger:
    def __init__(self):
        self.monitor = ProductionMonitor()
        self.registry = FLModelRegistry()
        self.last_retraining = None
    
    def should_retrain(self) -> tuple[bool, str]:
        """Determine if retraining is needed"""
        reasons = []
        
        # 1. Scheduled retraining (weekly)
        if self._is_scheduled_retraining():
            reasons.append("scheduled_weekly")
        
        # 2. Performance degradation
        if self.monitor.check_performance_degradation():
            reasons.append("performance_degradation")
        
        # 3. Data drift detected
        if self._check_data_drift():
            reasons.append("data_drift")
        
        # 4. Time since last retraining
        if self._time_based_retraining():
            reasons.append("time_based")
        
        # 5. New attack patterns (if you have this detection)
        if self._new_attack_patterns():
            reasons.append("new_attack_patterns")
        
        return len(reasons) > 0, ", ".join(reasons)
    
    def _is_scheduled_retraining(self) -> bool:
        """Check if it's time for scheduled retraining"""
        # Weekly retraining on Sundays
        return datetime.now().weekday() == 6  # Sunday
    
    def _check_data_drift(self) -> bool:
        """Check if data drift detected"""
        # Implementation would compare current data distribution
        # with baseline
        return False  # Placeholder
    
    def _time_based_retraining(self, days=7) -> bool:
        """Retrain if X days passed since last retraining"""
        if self.last_retraining is None:
            return True
        
        days_since = (datetime.now() - self.last_retraining).days
        return days_since >= days
    
    def _new_attack_patterns(self) -> bool:
        """Detect new attack patterns"""
        # This would require attack pattern analysis
        return False  # Placeholder
    
    def trigger_retraining(self, model_type: str):
        """Trigger FL retraining"""
        import subprocess
        
        # Start FL training
        subprocess.run([
            "python", "flower_server_metrics.py",
            "--model-type", model_type,
            "--num-rounds", "5"
        ])
        
        self.last_retraining = datetime.now()
\end{lstlisting}

\subsubsection{Step 6.2: Retraining Scheduler}

\textbf{File:} Create `mlops/retraining_scheduler.py`

\textbf{Code:}
\begin{lstlisting}[language=Python]
import schedule
import time
from mlops.retraining_trigger import RetrainingTrigger

class RetrainingScheduler:
    def __init__(self):
        self.trigger = RetrainingTrigger()
        self.model_types = ['MLPv2', 'LSTM', 'CNN1D', 'CNN_LSTM']
    
    def setup_schedule(self):
        """Setup retraining schedule"""
        # Weekly retraining on Sunday at 2 AM
        schedule.every().sunday.at("02:00").do(self.weekly_retraining)
        
        # Daily check for performance issues
        schedule.every().day.at("03:00").do(self.check_and_retrain_if_needed)
    
    def weekly_retraining(self):
        """Weekly scheduled retraining"""
        print("Starting weekly FL retraining...")
        for model_type in self.model_types:
            print(f"Retraining {model_type}...")
            self.trigger.trigger_retraining(model_type)
    
    def check_and_retrain_if_needed(self):
        """Check if retraining needed and trigger if so"""
        for model_type in self.model_types:
            should_retrain, reason = self.trigger.should_retrain()
            if should_retrain:
                print(f"Retraining {model_type}: {reason}")
                self.trigger.trigger_retraining(model_type)
    
    def run(self):
        """Run scheduler"""
        self.setup_schedule()
        while True:
            schedule.run_pending()
            time.sleep(60)  # Check every minute

if __name__ == "__main__":
    scheduler = RetrainingScheduler()
    scheduler.run()
\end{lstlisting}

\subsection{Phase 7: Model Comparison \& A/B Testing (Week 7)}

\subsubsection{Step 7.1: Model Comparison Tool}

\textbf{File:} Create `mlops/model_comparison.py`

\textbf{Code:}
\begin{lstlisting}[language=Python]
import mlflow
from mlflow.tracking import MlflowClient
import pandas as pd

class ModelComparator:
    def __init__(self):
        self.client = MlflowClient()
    
    def compare_models(self, model_type: str, num_versions: int = 5):
        """Compare recent model versions"""
        model_name = f"{model_type}_FL"
        
        # Get recent versions
        versions = self.client.search_model_versions(
            f"name='{model_name}'",
            max_results=num_versions
        )
        
        comparison = []
        for version in versions:
            run = self.client.get_run(version.run_id)
            comparison.append({
                "version": version.version,
                "round": run.data.params.get("round", "unknown"),
                "accuracy": run.data.metrics.get("accuracy", 0),
                "loss": run.data.metrics.get("loss", 0),
                "timestamp": version.creation_timestamp
            })
        
        df = pd.DataFrame(comparison)
        return df.sort_values("timestamp", ascending=False)
    
    def get_best_model(self, model_type: str):
        """Get best performing model version"""
        df = self.compare_models(model_type)
        if len(df) > 0:
            best = df.loc[df['accuracy'].idxmax()]
            return best
        return None
\end{lstlisting}

\subsubsection{Step 7.2: A/B Testing Framework}

\textbf{File:} Create `mlops/ab_testing.py`

\textbf{Code:}
\begin{lstlisting}[language=Python]
import random
from mlops.model_registry import FLModelRegistry

class ABTesting:
    def __init__(self, traffic_split: float = 0.5):
        """
        traffic_split: Fraction of traffic to new model (0.5 = 50/50 split)
        """
        self.traffic_split = traffic_split
        self.registry = FLModelRegistry()
        self.results = {
            "model_a": {"predictions": 0, "correct": 0},
            "model_b": {"predictions": 0, "correct": 0}
        }
    
    def route_request(self, model_type: str, features: List):
        """Route request to A or B model"""
        # Get current production model (A)
        model_a = self.registry.get_latest_model(model_type)
        
        # Get new candidate model (B)
        # This would be the newly trained model
        
        # Route based on traffic split
        use_new_model = random.random() < self.traffic_split
        
        if use_new_model:
            return "model_b", model_b
        else:
            return "model_a", model_a
    
    def record_result(self, model_version: str, prediction: int, 
                     actual: int, correct: bool):
        """Record A/B test result"""
        self.results[model_version]["predictions"] += 1
        if correct:
            self.results[model_version]["correct"] += 1
    
    def get_winner(self, confidence_level: float = 0.95) -> str:
        """Determine winner of A/B test"""
        from scipy import stats
        
        model_a_acc = (
            self.results["model_a"]["correct"] / 
            max(self.results["model_a"]["predictions"], 1)
        )
        model_b_acc = (
            self.results["model_b"]["correct"] / 
            max(self.results["model_b"]["predictions"], 1)
        )
        
        # Statistical significance test
        # (Simplified - would need proper statistical test)
        if model_b_acc > model_a_acc:
            return "model_b"
        return "model_a"
\end{lstlisting}

\section{Implementation Checklist}

\subsection{Phase 1: Foundation}
\begin{itemize}
    \item[ ] Install MLflow
    \item[ ] Set up MLflow tracking server
    \item[ ] Integrate MLflow into FL server
    \item[ ] Implement model versioning
    \item[ ] Set up structured logging
    \item[ ] Test experiment tracking
\end{itemize}

\subsection{Phase 2: Data Management}
\begin{itemize}
    \item[ ] Set up DVC for data versioning
    \item[ ] Implement data validation
    \item[ ] Create data quality monitoring
    \item[ ] Set up data drift detection
    \item[ ] Test data validation pipeline
\end{itemize}

\subsection{Phase 3: CI/CD}
\begin{itemize}
    \item[ ] Create GitHub Actions workflow
    \item[ ] Set up pre-commit hooks
    \item[ ] Write unit tests
    \item[ ] Configure automated testing
    \item[ ] Test CI/CD pipeline
\end{itemize}

\subsection{Phase 4: Deployment}
\begin{itemize}
    \item[ ] Create FastAPI serving endpoint
    \item[ ] Dockerize model serving
    \item[ ] Add to docker-compose
    \item[ ] Test API endpoints
    \item[ ] Deploy to staging environment
\end{itemize}

\subsection{Phase 5: Monitoring}
\begin{itemize}
    \item[ ] Implement production monitoring
    \item[ ] Set up alerting system
    \item[ ] Configure Prometheus metrics
    \item[ ] Create monitoring dashboard
    \item[ ] Test alerting
\end{itemize}

\subsection{Phase 6: Automation}
\begin{itemize}
    \item[ ] Implement retraining triggers
    \item[ ] Set up retraining scheduler
    \item[ ] Configure automated retraining
    \item[ ] Test retraining pipeline
\end{itemize}

\subsection{Phase 7: Advanced}
\begin{itemize}
    \item[ ] Implement model comparison
    \item[ ] Set up A/B testing
    \item[ ] Create model selection logic
    \item[ ] Test A/B testing framework
\end{itemize}

\section{File Structure After Implementation}

\begin{lstlisting}[language=text]
DDoS_SDN by Aiken Kazin/
├── mlops/
│   ├── __init__.py
│   ├── model_registry.py          # Model versioning
│   ├── fl_logger.py               # Structured logging
│   ├── data_validation.py         # Data quality checks
│   ├── data_quality.py            # Data drift detection
│   ├── model_serving.py           # FastAPI serving
│   ├── production_monitor.py      # Production monitoring
│   ├── alerting.py                # Alert system
│   ├── metrics_exporter.py        # Prometheus metrics
│   ├── retraining_trigger.py     # Retraining logic
│   ├── retraining_scheduler.py   # Scheduled retraining
│   ├── model_comparison.py        # Model comparison
│   └── ab_testing.py              # A/B testing
├── tests/
│   ├── test_fl_pipeline.py
│   ├── test_data_validation.py
│   └── test_model_serving.py
├── .github/
│   └── workflows/
│       └── fl_training.yml        # CI/CD pipeline
├── .pre-commit-config.yaml        # Pre-commit hooks
├── mlruns/                        # MLflow artifacts
├── mlflow.db                      # MLflow database
└── requirements-mlops.txt         # MLOps dependencies
\end{lstlisting}

\section{Success Metrics}

\subsection{Phase 1 Success Criteria}
\begin{itemize}
    \item All FL training runs logged to MLflow
    \item Models versioned and searchable
    \item Round-by-round metrics tracked
\end{itemize}

\subsection{Phase 2 Success Criteria}
\begin{itemize}
    \item Data validation catches issues before training
    \item Data drift detected automatically
    \item Dataset versions tracked
\end{itemize}

\subsection{Phase 3 Success Criteria}
\begin{itemize}
    \item CI/CD pipeline runs automatically
    \item Tests pass before deployment
    \item Automated model training on schedule
\end{itemize}

\subsection{Phase 4 Success Criteria}
\begin{itemize}
    \item Models accessible via API
    \item API response time < 100ms
    \item API uptime > 99\%
\end{itemize}

\subsection{Phase 5 Success Criteria}
\begin{itemize}
    \item Production metrics collected
    \item Alerts trigger on issues
    \item Monitoring dashboard functional
\end{itemize}

\subsection{Phase 6 Success Criteria}
\begin{itemize}
    \item Retraining triggers automatically
    \item Models retrain on schedule
    \item Performance improves over time
\end{itemize}

\section{Timeline Summary}

\begin{table}[h]
\centering
\begin{tabular}{|l|l|l|}
\hline
\textbf{Phase} & \textbf{Duration} & \textbf{Key Deliverables} \\
\hline
Phase 1: Foundation & Week 1-2 & MLflow integration, model versioning \\
Phase 2: Data Management & Week 2-3 & Data validation, versioning \\
Phase 3: CI/CD & Week 3-4 & Automated pipelines, testing \\
Phase 4: Deployment & Week 4-5 & Model serving API \\
Phase 5: Monitoring & Week 5-6 & Production monitoring, alerts \\
Phase 6: Automation & Week 6-7 & Automated retraining \\
Phase 7: Advanced & Week 7 & A/B testing, model comparison \\
\hline
\textbf{Total} & \textbf{7 weeks} & \textbf{Full MLOps pipeline} \\
\hline
\end{tabular}
\caption{Implementation Timeline}
\end{table}

\section{Next Steps}

\begin{enumerate}
    \item \textbf{Start with Phase 1:} Install MLflow and integrate into FL server
    \item \textbf{Test incrementally:} Test each phase before moving to next
    \item \textbf{Document as you go:} Update documentation with each change
    \item \textbf{Iterate:} Refine based on actual usage
\end{enumerate}

\end{document}

